\input hdr.tex
\section{A coupled partial melting Stokes-Darcy system}
Strong form
\begin{equation}
\begin{aligned}
		  \tilde{\bold{v}}_r + \frac{k_0 \phi_f^{1+\Theta}}{\mu_f} \nabla (\phi_f^{-1/2}\tilde{q}_f) &= 0 \\
		  \mu_s\phi_f^{-1/2}\nabla\cdot(\phi_f^{1+\Theta}\tilde{\bold{v}}_r) + \frac{1}{1-\phi_f}(\tilde{q}_f - \phi_f^{1/2}q) &=0\\
		  \nabla q - \nabla \cdot \hat{\sigma}(\bold{v}_s) &= -(1-\phi_f)\rho_r \bold{g} \\
		  \mu_s \nabla \cdot \bold{v}_s - \frac{\phi_f^{1/2}}{1-\phi_f}(\tilde{q}_f - \phi_f^{1/2}q) &=0
\end{aligned}
\end{equation}

List SI units

\begin{itemize}
		  \item [q] = pascal.
		  \item [$\bold{v}$] = m/s
\end{itemize}

Non dimensionalization choice
\begin{itemize}
		  \item $\check{\bold{v}} = \frac{\tilde{\bold{v}_r}}{k_0 \rho_r g_y/\mu_f}$
					 \item $\check{q} = \frac{q}{\rho_r g_y l_0 }$
					 \item $\check{x} = \frac{x}{l_0}$
					 \item $l_0 = (\frac{k_0 \mu_s}{\mu_f})^{1/2}$
\end{itemize}

The resulting non dimensionalized form is therefore
\begin{equation}
		  \begin{aligned}
					 \check{\bold{v}}_r + \phi_f^{1+\Theta}\nabla (\phi_f^{-1/2}\check{q}_f) &= 0\\
					 \phi_f^{-1/2}\nabla \cdot (\phi_f^{1+\Theta} \check{\bold{v}}_r) + \frac{1}{1-\phi_f}(\check{q}_f - \phi_f^{1/2}\check{q}) &=0 \\
					 \nabla \check{q} - \nabla \cdot \hat{\sigma}(\check{\bold{v}}_s) &= -(1-\phi_f)\rho_r \hat{\bold{g}} \\
					 \nabla \cdot \check{\bold{v}}_s - \frac{\phi_f^{1/2}}{1-\phi_f}(\check{q}_f - \phi_f^{1/2}\check{q}) &= 0
		  \end{aligned}
\end{equation}

scaled nondimensionalized weak form
let $\check{\sigma} = \mathcal{D}\bold{v} - \frac{1}{3}\nabla \cdot \bold{v} \bold{\text{I}}$

\begin{equation}
		  \begin{aligned}
					 (2\phi_s\check{\sigma}(\bold{v}_{s,h}),\nabla \Psi_s) - (\check{q}_h, \nabla \cdot \Psi_s) &= -(\phi_s\check{\bold{f}}, \Psi_s) \\
					 (\nabla \cdot \check{\bold{v}}_{s,h}, \omega)  + \Big(\frac{\hat{\phi}_f}{\phi_s} \check{q}_h, \omega\Big) - 
					\Big(\frac{\hat{\phi}_f^{1/2}}{\phi_s} \check{q}_{f,h},\omega \Big) &= 0\\
					 (\bold{\check{v}}_{r,h},\Psi_r)  - (\check{q}_{f,h},\hat{\phi}_f^{-1/2} \nabla \cdot (\phi_f^{1+\Theta}\Psi_r)) &= 0\\
					 (\hat{\phi}_f^{-1/2} \nabla \cdot (\phi_f^{1+\Theta}\Psi_r),\omega_f) + \Big(\frac{1}{\phi_s}\check{q}_{f,h},\omega_f \Big) - \Big( \frac{\hat{\phi}_f^{1/2}}{\phi_s} \check{q}_h, \omega_f \Big) &= 0
		  \end{aligned}
\end{equation}

\newpage

\section{Linear system and Uzawa type algorithm}

original linear system
\begin{equation}
\begin{bmatrix}
\bold{A}_s   & -\bold{B}_s & \bold{0}     &  \bold{0}   \\
\bold{B}_s^T &  \bold{C}_s & \bold{0}     & -\bold{K}   \\
\bold{0}     &  \bold{0}   & \bold{A}_d   & -\bold{B}_d \\
\bold{0}     & -\bold{K}   & \bold{B}_d^T &  \bold{C}_d
\end{bmatrix} \begin{bmatrix}
\bold{x}_s\\
\bold{y}_s\\
\bold{x}_d\\
\bold{y}_d
\end{bmatrix}= 
\begin{bmatrix}
\bold{f}\\
\bold{0}\\
\bold{0}\\
\bold{0}
\end{bmatrix}
\end{equation}

This is a double saddle point system. 
Symmetrically permutate original linear system and form a enlarged saddle point system.
\begin{equation}
\begin{bmatrix}
\bold{A}_s   & \bold{0}     & -\bold{B}_s &  \bold{0}   \\
\bold{0}     & \bold{A}_d   &  \bold{0}   & -\bold{B}_d \\
\bold{B}_s^T & \bold{0}     &  \bold{C}_s & -\bold{K}   \\
\bold{0}     & \bold{B}_d^T & -\bold{K}   &  \bold{C}_d
\end{bmatrix} \begin{bmatrix}
\bold{x}_s\\
\bold{x}_d\\
\bold{y}_s\\
\bold{y}_d
\end{bmatrix}= 
\begin{bmatrix}
\bold{f}\\
\bold{0}\\
\bold{0}\\
\bold{0}
\end{bmatrix}
\end{equation}

\begin{equation}
\tilde{\bold{A}} = \begin{bmatrix}
\bold{A}_s   & \bold{0}    \\
\bold{0}     & \bold{A}_d  
\end{bmatrix}
\quad
\tilde{\bold{B}} = \begin{bmatrix}
\bold{B}_s & \bold{0}   \\
\bold{0}   & \bold{B}_d
\end{bmatrix}
\quad
\tilde{\bold{C}} = \begin{bmatrix}
 \bold{C}_s & -\bold{K}   \\
-\bold{K}   &  \bold{C}_d
\end{bmatrix}
\end{equation}
In this newly formed system, the $\tilde{\bold{A}}$ matrix maintains as a positive definite matrix.
we solve this typical saddle point system with uzawa iteration. Uzawa algorithm has been discussed
extensively at \cite{InexactUzawa_Elam_Golub}, \cite{AnalysisInexactUzawa_Bramble_Pasciak_Vassilev}.

Consider the linear system of equation:
\begin{equation}
\begin{bmatrix}
\bold{A}   & -\bold{B} \\
\bold{B}^T &  \bold{C}
\end{bmatrix} \begin{bmatrix}
\bold{x} \\ \bold{y}
\end{bmatrix} = \begin{bmatrix}
\bold{f} \\
\bold{0}
\end{bmatrix}
\end{equation}

\begin{algorithm}
\caption{Inexact Uzawa Iteration}
Initialize solution
\end{algorithm}

\subsection{Convergence and stability analysis}
Uzawa iteration is in fact a fixed point iteration procedure.
Usually, uzawa iteration convergences slowly.

\subsubsection{Anderson acceleration}
Accelerating uzawa iteration has been discussed in the previously literatures.
Anderson acceleration (Anderson mixing) \cite{AcceleratingUzawa_ho_Olson_Walker}

\newpage

\section{Test cases}
We test the linear system a little bit different from the linear system above. 
\subsection{Constant porosity and balanced pressure (divergence free velocity field)}


\subsection{Const porosity and unbalanced pressure}
\begin{equation}
\check{q} = 0 \quad \quad \check{q}_f = 2(x+y) \quad \quad 
\check{\bold{v}}_r = -\frac{\phi_f^{-1/2}}{1-\phi_f}\begin{pmatrix}x^2 \\ y^2 \end{pmatrix} 
\quad \quad \check{\bold{v}}_s = \frac{\phi_f^{1/2}}{1-\phi_f} \begin{pmatrix}x^2 \\ y^2 \end{pmatrix}
\end{equation}
\begin{equation}
\nabla \cdot \check{\bold{v}}_r = \frac{-2\phi_f^{-1/2}}{1-\phi_f}(x+y) 
\end{equation}
\begin{equation}
\mathcal{D}\check{\bold{v}}_s = \frac{2\phi_f^{1/2}}{1-\phi_f}\begin{bmatrix} x & 0 \\ 0 & y \end{bmatrix}
\quad \quad 
(\nabla \cdot \check{\bold{v}}_s) \bold{\text{I}} = \frac{2\phi_f^{1/2}}{1-\phi_f} \begin{bmatrix} x+y & 0 \\ 0 & x+y \end{bmatrix}
\end{equation}

\begin{equation}
\begin{aligned}
\tilde{\bold{v}}_r + \phi_f^{1/2}\nabla\check{q}_f &= 
-\frac{\phi_f^{-1/2}}{1-\phi_f}\begin{pmatrix}x^2 \\ y^2 \end{pmatrix} + \phi_f^{1/2}\begin{pmatrix} 2 \\ 2 \end{pmatrix}\\
\phi_f^{1/2}\nabla \cdot (\check{\bold{v}}_r) + \frac{1}{1-\phi_f}\check{q}_f  &= \frac{-2}{1-\phi_f}(x+y) + \frac{1}{1-\phi_f}2(x+y)=0\\
\nabla \cdot \check{\bold{v}}_s - \frac{\phi_f^{1/2}}{1-\phi_f}\check{q}_f &= \frac{2\phi_f^{1/2}}{1-\phi_f}(x+y) - \frac{2\phi_f^{1/2}}{1-\phi_f}(x+y) = 0\\
-2(1-\phi_f)\nabla \cdot \sigma(\check{\bold{v}}_s) &= -2(1-\phi_f)\nabla \cdot (\mathcal{D}\check{\bold{v}}_s - \frac{1}{3}(\nabla \cdot \check{\bold{v}}_s) \bold{\text{I}}) = 
-\frac{8\phi_f^{1/2}}{3} \begin{pmatrix} 1 \\ 1 \end{pmatrix}
\end{aligned}
\end{equation}

% ====================
\newpage
\bibliographystyle{siam}
\bibliography{ref}

\end{document}
